\documentclass{article}
\usepackage[utf8]{inputenc}
\usepackage{xcolor}
\usepackage{hyperref}

\usepackage{rotating}

\usepackage{tabls}
\usepackage{booktabs}
\usepackage{amsmath}
\usepackage{amssymb}
\usepackage{amsthm}
\usepackage{amsfonts}
\usepackage{multicol}
\usepackage{enumitem}
\usepackage{microtype}
\usepackage{tikz}
\usepackage{pgfplots}
\usetikzlibrary{positioning}
\usepackage{multirow}

\usepackage{comment}

\usepackage{graphicx}
\usepackage{wrapfig}

\theoremstyle{plain}
\newtheorem{theorem}{Theorem}
\newtheorem*{theorem*}{Theorem}
\newtheorem{lemma}[theorem]{Lemma}
\newtheorem*{lemma*}{Lemma}
\newtheorem{proposition}[theorem]{Proposition}
\newtheorem{corollary}[theorem]{Corollary}

\theoremstyle{box}
\newtheorem{definition}[theorem]{Definition}
\newtheorem*{axiom*}{Axiom}
\newtheorem{dfn}[theorem]{Definition}
\newtheorem*{definition*}{Definition}
\newtheorem*{exercise*}{Exercise}
\newtheorem{observation}[theorem]{Observation}
\newtheorem{remark}[theorem]{Remark}
\newtheorem{example}[theorem]{Example}
\newtheorem{question}[theorem]{Question}
\newtheorem*{prep-problems}{Preparation Problems}

\newcommand{\pd}[3]{\frac{\partial^{#3} {#1}}{\partial {#2}^{#3}}}
\newcommand{\fpd}[2]{\frac{\partial {#1}}{\partial {#2}}}
\newcommand{\diff}[2]{\frac{d {#1}}{d {#2}}}


\begin{document}
The Project 1 template includes background for the data and questions we will be exploring. To complete this project you will need to complete the following tasks insuring that you meet \textbf{ALL} the specifications. If \textbf{any} of the specification are not met you will receive feedback identifying the specification(s) that still require work in order to complete the project.

\begin{itemize}
\item Tell a Story with Functions
	\begin{itemize}
	\item Task 1: Create a scatterplot of the data for one light bulb. \textcolor{red}{(which light bulb can be randomly generated for each student)}
	\item Task 2a: Plot the five deterministic models.
		\begin{itemize}
		\item $f_1(t) = 100$
		\item $f_2(t) = 100 - 0.003t$
		\item $f_3(t) = 120 - 20e^{-0.0004t}$
		\item $f_4(t) = 95 - 0.002t + 2\ln(18t+100)$
		\item $f_5(t) = (100+0.007t)e^{-0.00005t}$
		\end{itemize}
	\item Task 2b: Describe in words the story told by each function in the context of the light bulbs.
	\end{itemize}
\item Solving Equations (with Software)
	\begin{itemize}
	\item Task 3a: Use R to solve each five deterministic functions, $f_i(t) = 75$.
	\item Task 3b: Identify which of the numeric results you can verify analytically.
	\item Task 3c: Explain your results in the context of the light bulbs.
	\end{itemize}
\item Modeling Measurement
	\begin{itemize}
	\item Task 4a: Plot a histogram of the data for all the light bulbs at a given time. \textcolor{red}{(which time can be randomly generated for each student)}
	\item Task 4b: Select values for $\mu$ and $\sigma$ to obtain a visual fit of the normal distribution, $f(L) = \frac{1}{\sigma\sqrt{2\pi}}e^{-\frac{1}{2}\left(\frac{L-\mu}{\sigma}\right)^2}$, to the histogram of the data.
	\item Task 4c: Create plots with both the data (the histogram) and a selected model (normal distribution with the parameter values you picked). Include your best visual fit and one poor visual fit. Make sure you clearly indicate which parameter values you used in each case.
	\item Task 4d: Explain what your best visual fit parameters mean in the context of the light bulb question.
	\end{itemize}
\item Selecting Parameter Values (for a Visual Fit) \\ for the model $f_5(t) = (a + bt)e^{-ct}$
	\begin{itemize}
	\item Task 5a: There are three parameter in this model -- $a$, $b$, and $c$. Since we know $f_5(0) = 100$ this tells us the value for one of these parameters. Identify which parameter we know and give its value.
	\item Task 5b: Substitute in the parameter value we know and then experiment with values for the remaining two parameters. Explain what you learned in your exploration. Include 4-6 key plots (constructed in R) to support your summary. \\
	Some questions you might consider to help you get started with your exploration include:
		\begin{itemize}
		\item How will you determine what window to use when making your graphs?
		\item Can you find any combinations of parameters that result in $f_5$ only decreasing during the time interval of interest? Can you find any combinations of parameters that result in $f_5$ only increasing during the time interval of interest? Can you find any combinations of parameters that result in $f_5$ increasing and then decreasing during the time interval of interest? Can you find any combinations of parameters that result in $f_5$ decreasing and then increasing during the time interval of interest? What do any of these results mean in the context of the light bulb story?
		\item Can you find any values of $b$ where the behavior/characteristics of $f_5$ are different than for other values of $b$? Can you find any values of $c$ where the behavior/characteristics of $f_5$ are different than for other values of $c$?
		\item Did you find any values of $b$ or $c$ (or combinations of values of $b$ and $c$) that don't make sense in the context of the light bulb story?
		\item Can you explain any of your observations using what you know about the five building block functions and 5 function operations that were used to build up the function $f_5$?
		\item What did you notice about the values of $b$ and $c$ (and their combinations) that you can used to help fine tune the behavior of $f_5$ to fit the context of the light bulb story?
		\end{itemize}
	\item Task 5c: Create plots with both the data (the scatterplot) and a selected model ($f_5$ with the parameter values you picked). Include your best visual fit and one poor visual fit. Make sure you clearly indicate which parameter values you used in each case.
	\item Task 5d: Explain what your best visual fit parameters mean in the context of the light bulb question.
	\end{itemize}
\item Conclusions
	\begin{itemize}
	\item Task 6: Answer the question, ``Do these light bulbs appear to meet the requirements for the L Prize?" Make sure to provide support for your answer using the models you selected and your data. 
	\end{itemize}
\item Reflection
	\begin{itemize}
	\item Task 7: Reflect on your work for this project. In 1-2 paragraphs, identify/explain 2-3 key mathematical  ideas you learned (and would like to remember) and 1-3 soft skills you needed/improved/learned while working on the project.
	\end{itemize}
\end{itemize}


 
\end{document}